\apartado{}{Fórmulas}{preambulo}
\bajada{Están acá y no al principio de cada apartado a propósito: si la fórmula
está en el renglón de arriba, el ejercicio se resuelve copiando. Ir a buscarla
obliga a decidir cuál hace falta, que es la mitad del trabajo.}

\definicion{Probabilidad}{%
  \textbf{Clásica} (sólo con simetría):
  $P(A)=\dfrac{\text{casos favorables}}{\text{casos posibles}}$\par
  \textbf{Frecuentista}:
  $P(A)\approx\dfrac{\text{veces que ocurrió}}{\text{veces que se observó}}$\par
  \textbf{Complemento}: $P(A^{c})=1-P(A)$\par
  \textbf{Axiomas}: $0\le P(A)\le 1$ \quad y \quad $P(S)=1$}

\definicion{Espacio muestral}{%
  Si un procedimiento tiene $k$ etapas y cada una tiene $m$ resultados
  posibles, el total es $m^{k}$ — se multiplica, no se suma.\par
  Ejemplo: 9 ítems de 4 opciones dan $4^{9}=262\,144$ formas de responder,
  y sólo 28 puntajes distintos.}

\definicion{Tabla de contingencia}{%
  \begin{center}\small
  \begin{tabular}{@{}lcc@{}}
  \toprule
   & \textbf{Condición sí} & \textbf{Condición no}\\
  \midrule
  \textbf{Test $+$} & VP & FP\\
  \textbf{Test $-$} & FN & VN\\
  \bottomrule
  \end{tabular}
  \end{center}
  $\text{Sensibilidad}=\dfrac{VP}{VP+FN}$ \quad
  $\text{Especificidad}=\dfrac{VN}{VN+FP}$ \quad
  $\text{VPP}=\dfrac{VP}{VP+FP}$\par
  Las tres tienen el mismo tipo de numerador y distinto denominador: las dos
  primeras dividen por \emph{columna}, el VPP por \emph{fila}.}

\definicion{Acuerdo entre evaluadores}{%
  $P_{o}$ = proporción de casos en que coinciden.\par
  $P_{e}$ = acuerdo esperado por azar = (proporción de «sí» de A $\times$
  proporción de «sí» de B) $+$ (proporción de «no» de A $\times$ proporción de
  «no» de B).\par
  $\kappa=\dfrac{P_{o}-P_{e}}{1-P_{e}}$}

\definicion{Combinatoria}{%
  $n! = n\times(n-1)\times\cdots\times 2\times 1$, \; y \; $0!=1$\par
  \textbf{Permutación} (importa el orden): $P(n,k)=\dfrac{n!}{(n-k)!}$\par
  \textbf{Combinación} (no importa): $C(n,k)=\dfrac{n!}{k!\,(n-k)!}$\par
  Relación: $P(n,k)=C(n,k)\times k!$}
